% !TEX root = ../../main.tex
% C5.3 - Database design (translated). Matches storage/postgres models and migrations (12 tables), storage/milvus/vectors.py,
% storage/minio/frames.py, services/track_ingestion.py and storage/maintenance_cli.py. Figures H6, H8 are the English versions.
\section{Database Design}
\label{sec:5.3}

\subsection{Relational Database}
\label{sec:5.3.1}

PostgreSQL stores all the business data and is the source of truth for permissions and states. The database has 12 tables in five groups, as in Table~\ref{tab:db-tables}.

\begin{table}[H]
\centering
\caption{Tables of the relational database}
\label{tab:db-tables}
\begin{tabularx}{\textwidth}{|>{\raggedright\arraybackslash}p{3.0cm}|>{\raggedright\arraybackslash}p{5.4cm}|L|}
\hline
\textbf{Group} & \textbf{Tables} & \textbf{Role} \\
\hline
Organisation and access & \texttt{areas}, \texttt{users}, \path{auth_sessions} & Surveillance areas, user accounts, login sessions \\
\hline
Cameras and AI processing & \texttt{cameras}, \path{ai_config_versions}, \path{processing_jobs}, \path{worker_heartbeats} & Logical cameras, versions of the Detector-Tracker-encoder configuration, job queue, background worker status \\
\hline
Search data & \path{person_tracks}, \path{storage_outbox_events} & Descriptions of the appearances; synchronisation events for the other two stores (Section~\ref{sec:5.3.4}) \\
\hline
Cases & \texttt{cases}, \path{case_results} & Cases and their saved results \\
\hline
Traceability & \path{audit_logs} & Audit log, append-only \\
\hline
\end{tabularx}
\end{table}

Figure~\ref{fig:erd} shows the entity-relationship diagram of the 11 tables that have foreign keys, in crow's foot notation; each table only shows its primary key, foreign keys and main attributes. The \path{worker_heartbeats} table has no foreign key and is not drawn.

\afterpage{%
\begin{landscape}
\begin{figure}[p]
\centering
\includegraphics[width=\linewidth,height=0.86\textheight,keepaspectratio]{H8-erd.png}
\caption{Entity-relationship diagram of the relational database}
\label{fig:erd}
\end{figure}
\end{landscape}
}

\paragraph{Relationships.} Most foreign keys are mandatory: each camera belongs to exactly one area, each case has exactly one owner, and each appearance belongs to exactly one camera and was created by exactly one job with one AI configuration version. Three foreign keys may be null: an account only has an area when it is an Operator, a job only has a requester when it was created from an uploaded file, and a log entry may not be tied to a user, for example when it records a technical error. \texttt{cases} and \path{person_tracks} have a many-to-many relationship through \path{case_results}. Each row of this table has its own primary key, stores a snapshot of the information at the time of saving, and has \emph{no} uniqueness constraint on the case-appearance pair, since each save creates a new entry.

\paragraph{Integrity constraints.} Many business rules are expressed directly in the database, so that invalid data cannot be written even if the application layer has a bug:
\begin{itemize}
    \item An Operator account always has an area and the other two roles never do; a case has a closing time if and only if it is \texttt{CLOSED}; only one AI configuration version is \texttt{ACTIVE}; the area code and the area of a camera do not change after creation.
    \item Most foreign keys use the RESTRICT rule, so data that is still referenced cannot be deleted; accounts and cameras are not deleted but change status (\texttt{INACTIVE}, \texttt{RETIRED}). Only login sessions and result entries are deleted with their parent.
    \item An appearance can only be \texttt{READY} when it has the object key, hash and size of its frame and the time its vector was written. No table stores a relevance score.
    \item The \texttt{users}, \texttt{cameras} and \texttt{cases} tables have a version column to prevent overwriting on concurrent updates (Section~\ref{sec:5.2.4}); the requester-idempotency key pair of \path{processing_jobs} is unique, so a repeated upload request creates no additional job.
\end{itemize}
Indexes are also created for frequent queries, such as appearances by camera and time (with a separate index containing only ready appearances), cases by owner, and log entries by actor.

\subsection{Feature Vector Storage}
\label{sec:5.3.2}

Milvus serves a single purpose: quickly finding the vectors closest to the query vector within a filtered scope; all other information and every permission decision live in PostgreSQL.

\paragraph{One collection per encoder version.} Milvus organises data into \emph{collections}, similar to tables. The system creates a separate collection for each encoder version, for example \path{person_track_embeddings_rasa_cuhk_pedes_v1}, because vectors from two versions belong to two representation spaces that cannot be compared. A fixed alias always points to the collection of the version in use; when the background worker starts, the dimension of the collection is checked against the configuration. Each appearance corresponds to one record with the fields of Table~\ref{tab:milvus-schema}.

\begin{table}[H]
\centering
\caption{Schema of the Milvus collection of feature vectors}
\label{tab:milvus-schema}
\begin{tabularx}{\textwidth}{|>{\raggedright\arraybackslash}p{4.2cm}|>{\raggedright\arraybackslash}p{3.7cm}|L|}
\hline
\textbf{Field} & \textbf{Type} & \textbf{Meaning} \\
\hline
\path{track_id} & VARCHAR(36), primary key & Appearance identifier, the same as the primary key of \path{person_tracks} \\
\hline
\path{embedding} & FLOAT\_VECTOR, 256 dimensions & Feature vector, L2-normalised \\
\hline
\path{area_id}, \path{camera_id} & VARCHAR(36) & Area and camera, used for scope filtering \\
\hline
\path{appeared_at_epoch} & INT64 & Time of appearance (Unix seconds), used for time filtering \\
\hline
\path{encoder_version} & VARCHAR(64) & Version of the encoder that produced the vector \\
\hline
\path{index_status} & VARCHAR(16) & Readiness; only \texttt{READY} records are searched \\
\hline
\end{tabularx}
\end{table}

The area, camera and time fields are copied from PostgreSQL so that Milvus can pre-filter before searching (Section~\ref{sec:3.6.3}); since the area of a camera does not change after creation, this copy cannot diverge. The operational status of a camera, by contrast, changes with the Administrator's actions and is not copied: each search takes the list of operational cameras from PostgreSQL and puts it in the filter, so the data of a retired camera is hidden immediately without updating Milvus.

\paragraph{Index and operations.} The vector field has an HNSW index with the inner-product metric, $M = 16$, \textit{ef\-Con\-struc\-tion}~$= 128$, and $\mathit{ef} = 64$ at search time (Section~\ref{sec:3.6.4}). Vectors are written with an \emph{upsert} keyed by the appearance identifier, so retries create no duplicate records, and are checked for the right dimension, the absence of invalid values and unit length before writing or searching. Every operation uses strong consistency, so a vector just written appears in the very next search.

\subsection{Frame Storage}
\label{sec:5.3.3}

Each appearance has exactly one image file in MinIO: the representative full frame, stored as JPEG at quality 90; PostgreSQL only stores the object key, hash and file size. The system does not store cropped person images as separate files: the person image in the result list is rendered from the full frame and the bounding box on request (Section~\ref{sec:6.3.3}). One file therefore serves both views, the cropping can change without regenerating data, and no standalone person images exist that could be copied out of control. The original video is not stored.

The key of each file is derived deterministically from the appearance, following the pattern
\begin{center}
\texttt{tracks/v1/}\allowbreak \textit{camera code}\texttt{/}\allowbreak \textit{year}\texttt{/}\allowbreak \textit{month}\texttt{/}\allowbreak \textit{day}\texttt{/}\allowbreak \textit{appearance id}\texttt{/}\allowbreak \texttt{representative.jpg}
\end{center}
so a retry rewrites the same file instead of creating a new one; the versioned prefix allows the layout to change later, and grouping by camera and day makes clean-up by camera or time range easy.

Before writing, the file is checked for the right format, a size of at most 20~MB and dimensions matching its description; its SHA-256 hash is stored in both MinIO and PostgreSQL. A write to an existing key with the same hash reuses the existing file, and one with a different hash is rejected as a conflict; on reading, the hash is recomputed to detect corruption. The files live in a private bucket, the application accesses it with an account that only has rights on that bucket, and the browser never receives a direct address to MinIO.

\subsection{Consistency Across the Stores}
\label{sec:5.3.4}

An appearance is stored as three parts in three stores, linked by the same identifier. The three stores share no transaction, so the process may stop midway or a store may be temporarily unreachable. The system guarantees that incomplete data never appears in search results and can be completed or detected later.

\paragraph{Appearance states.} Each appearance record has a synchronisation state, with the allowed transitions of Figure~\ref{fig:track-state}. Only appearances in the \emph{Ready} state can be searched; transitions not shown in the figure are rejected by the application layer.

\begin{figure}[!htbp]
\centering
\includegraphics[width=\textwidth]{H6-trang-thai-lan-xuat-hien.png}
\caption{State diagram of an appearance}
\label{fig:track-state}
\end{figure}

\paragraph{Write order.} An appearance is written in a fixed order:
\begin{enumerate}
    \item In \emph{one} PostgreSQL transaction, create the record in the \emph{Pending} state and an event in \path{storage_outbox_events} holding everything that must be written to the other two stores, the vector included. This is the \emph{transactional outbox} pattern: the intent to write to another system is stored in the same transaction as the main data.
    \item Write the frame to MinIO under its deterministic key.
    \item Write the vector to Milvus, then read it back to confirm the record exists with the right area and camera.
    \item In a new transaction, move the record to \emph{Ready}, store the hash and size of the frame and the time the vector was written, and mark the outbox event as done.
\end{enumerate}
Because PostgreSQL is always written \emph{first}, every frame and vector can be matched with a known record, and anything without a matching record is orphaned data. Because the \emph{Ready} state is set \emph{last}, an appearance can only be found once all its data is present. Steps 2 and 3 can both be repeated safely thanks to the deterministic key and the upsert.

\paragraph{Error handling and maintenance tools.} If step 2 or 3 fails with a transient error, the record stays \emph{Pending} and the outbox event is scheduled for retry, with a delay starting at 2 seconds, doubling each time, up to 300 seconds. After five unsuccessful attempts, or on an unrecoverable error such as a hash conflict, the record moves to \emph{Failed} and the incident is written to the audit log. Since the outbox event holds all the data, a retry does not need to rerun the video analysis. Operators also have a command-line tool to complete pending appearances; to reconcile the three stores, finding records pending for too long, ready records missing their frame or vector, and orphaned data (by default it only reports; on request it quarantines broken records and deletes orphans); to move a failed appearance back to pending; and to rebuild the Milvus data from the outbox events. These commands are currently run by hand.

\paragraph{Protection at search time.} Even when the three stores are temporarily out of sync, users never see wrong data: every result returned by Milvus is checked again against PostgreSQL, and results that do not exist or are outside the permitted scope are removed before being shown (Section~\ref{sec:3.6.4}).
